%% toc: P2 10 10.1,10.2
\section{Proofs for readability and valuations over an arbitrary set of prime formulas}
\label{sec:annex-p02-w02b-p02-001-prime-readability}

The notation is that of \cref{sec:p02-w02b-p02-001-prime-readability}.  For
an expression \(X\), \(|X|\) is its length and \(X[k]\) its \(k\)-th entry.
An expression \(E\) \emph{occurs in \(X\) at position \(k\)} if
\(X[k+i-1]=E[i]\) for \(1\leq i\leq|E|\).  The following observation is used
in every proof below.

\emph{Observation (O).}  Let \(X\) be an initial segment of \(X'\), and let
\(E\) occur in \(X\) at position \(k\) and \(E'\) occur in \(X'\) at the same
position \(k\).  Then the shorter of \(E,E'\) is an initial segment of the
longer.  Indeed, if \(|E|\leq|E'|\), the entries
\(X[k],\ldots,X[k+|E|-1]\) of \(E\) lie in \(X\), where \(X\) and \(X'\)
agree, so they are the first \(|E|\) entries of \(E'\).  If \(|E'|<|E|\),
the entries of \(E'\) occupy positions up to \(k+|E'|-1<k+|E|-1\leq|X|\),
again inside \(X\), so \(E'\) consists of the first \(|E'|\) entries of
\(E\).

\begin{annexitem}{annex:p02-w02b-p02-001-readability}{thm:p02-w02b-p02-001-letter-readability}{Unique readability}
\begin{proof}
Write \(F\) for the set of propositional formulas of \(P\).

\emph{Step 1: inversion.}  Let \(F^{*}\) be the set of \(A\in F\) that are
a member of \(P\), or are \(\neg B\) with \(B\in F\), or are \((B\circ D)\)
with \(B,D\in F\) and \(\circ\) a binary connective.  Every member of \(P\)
lies in \(F^{*}\).  If \(A,B\in F^{*}\) then \(A,B\in F\), so \(\neg A\) and
\((A\circ B)\) lie in \(F\) and are of the second and third kinds; hence
they lie in \(F^{*}\).  Thus \(F^{*}\) contains \(P\) and is closed under the
formation rules, and since \(F\) is the smallest such set, \(F\subseteq
F^{*}\); so \(F=F^{*}\).  Consequently every formula is nonempty, and its
first entry is a member of \(P\), the symbol \(\neg\), or a left
parenthesis, according as it is of the first, second or third kind.  These
three possibilities exclude one another because \(P\) contains none of the
six symbols.  In particular a formula whose first entry is \(p\in P\) is
the expression \(p\) of length \(1\).

\emph{Step 2: no formula is a proper initial segment of another.}  We
prove by strong induction on \(|A|\): for every \(A'\in F\), if \(A\) is an
initial segment of \(A'\) then \(A=A'\).  Let \(A\) be an initial segment
of \(A'\), so \(|A|\leq|A'|\) and the two share their first entry.

If the first entry is \(p\in P\), then \(A=p=A'\) by Step 1.

If the first entry is \(\neg\), then \(A=\neg B\) and \(A'=\neg B'\) with
\(B,B'\in F\), occurring at position \(2\).  By (O) the shorter of \(B,B'\)
is an initial segment of the longer.  The shorter has length below \(|A|\)
(as \(|B|<|A|\) and the shorter is at most \(|B|\)), so the induction
hypothesis applied to it gives \(B=B'\), and \(A=A'\).

If the first entry is a left parenthesis, then \(A=(B\circ D)\) and
\(A'=(B'\circ' D')\).  The formulas \(B\) and \(B'\) occur at position \(2\);
by (O) the shorter is an initial segment of the longer, and it is shorter
than \(A\), so the induction hypothesis gives \(B=B'\).  The entry at
position \(|B|+2\) is \(\circ\) in \(A\) and \(\circ'\) in \(A'\); this
position is at most \(|A|\), so \(\circ=\circ'\).  Now \(D\) and \(D'\)
occur at position \(|B|+3\); by (O) and the induction hypothesis applied to
the shorter of the two, which is shorter than \(A\), \(D=D'\).  Hence
\(A=A'\).

\emph{Step 3: readability.}  By Step 1 each formula is of at least one of
the three kinds, and the first entry shows that it is of exactly one.  If
\(\neg B=\neg B'\), deleting the first entry gives \(B=B'\).  If
\((B\circ D)=(B'\circ'D')\), then \(B\) and \(B'\) occur at position \(2\)
of the same expression, so by (O) one is an initial segment of the other,
and Step 2 gives \(B=B'\).  Then \(\circ\) and \(\circ'\) are the entry at
position \(|B|+2\), and \(D\) and \(D'\) both consist of the entries from
position \(|B|+3\) to the second-last, so \(\circ=\circ'\) and \(D=D'\).
Step 2 is the last assertion of the theorem.
\end{proof}
\end{annexitem}

\begin{annexitem}{annex:p02-w02b-p02-001-fo-read}{thm:p02-w02b-p02-001-fo-skeleton}{First-order readability}
\begin{proof}
Call a symbol a \emph{term symbol} if it is a variable, a constant symbol
or a function symbol of \(L\).  By hypothesis the classes of variables,
constant symbols, function symbols and relation symbols are disjoint from
one another and from the logical symbols and the parentheses, and each
function or relation symbol has a single arity.

\emph{Inversion for terms.}  By the argument of Step 1 of
\cref{annex:p02-w02b-p02-001-readability}, applied to the smallest
set of terms in \cref{def:p02-w02-p02-005-bw-syntax}, every term is a variable, a constant symbol, or
\(f(t_1\cdots t_n)\) with \(f\) an \(n\)-ary function symbol and
\(t_1,\ldots,t_n\) terms: the terms of these three shapes contain the
variables and constants and are closed under the formation rule.  A term
whose first symbol is a variable or a constant symbol is that one symbol,
since a term of the third shape begins with a function symbol.

\emph{Claim T: if a term \(t\) is an initial segment of a term \(t'\), then
\(t=t'\).}  Strong induction on \(|t|\).  The first symbols agree.  If it
is a variable or constant symbol, \(t\) and \(t'\) are both that symbol.
Otherwise it is a function symbol \(f\), of arity \(n\), and
\(t=f(t_1\cdots t_n)\), \(t'=f(t'_1\cdots t'_n)\).  We show
\(t_i=t'_i\) for \(i=1,\ldots,n\) in turn.  If \(t_j=t'_j\) for all
\(j<i\), then \(t_i\) and \(t'_i\) occur at the same position
\(3+|t_1|+\cdots+|t_{i-1}|\) of \(t\) and \(t'\).  By (O) the shorter is an
initial segment of the longer; it is shorter than \(t\), so the induction
hypothesis gives \(t_i=t'_i\).  With all arguments equal, \(t=t'\).

\emph{Inversion for formulas.}  In the same way, every formula of \(L\) is
of one of the forms \((t_1=t_2)\), \(R(t_1\cdots t_n)\), \(\neg\varphi\),
\((\varphi\circ\psi)\), \((\exists v\varphi)\), \((\forall v\varphi)\), with
\(t_i\) terms, \(R\) an \(n\)-ary relation symbol, \(\varphi,\psi\)
formulas and \(v\) a variable.  Hence the first symbol of a formula is a
left parenthesis, a relation symbol or \(\neg\).  For the forms beginning
with a left parenthesis, the second symbol is a term symbol for
\((t_1=t_2)\) (the first symbol of \(t_1\), by inversion for terms), is
\(\exists\) or \(\forall\) for the quantified forms, and is a left
parenthesis, a relation symbol or \(\neg\) for \((\varphi\circ\psi)\) (the
first symbol of \(\varphi\)).  These three sets of second symbols are
pairwise disjoint.  So the first two symbols of a formula determine which
of the six forms it has.

\emph{Claim F: if a formula \(\varphi\) is an initial segment of a formula
\(\varphi'\), then \(\varphi=\varphi'\).}  Strong induction on
\(|\varphi|\).  Every formula has length at least \(2\), so \(\varphi\) and
\(\varphi'\) share their first two symbols and hence their form.
\begin{itemize}
\item \(R(t_1\cdots t_n)\) and \(R(t'_1\cdots t'_n)\): as in Claim T, the
  arguments occur at matching positions in turn, and (O) with Claim T gives
  \(t_i=t'_i\) for each \(i\).
\item \((t_1=t_2)\) and \((t'_1=t'_2)\): \(t_1,t'_1\) occur at position
  \(2\), so (O) and Claim T give \(t_1=t'_1\); then \(t_2,t'_2\) occur at
  position \(|t_1|+3\), and the same argument gives \(t_2=t'_2\).
\item \(\neg\chi\) and \(\neg\chi'\): \(\chi,\chi'\) occur at position \(2\);
  (O) and the induction hypothesis, applied to the shorter of the two,
  which is shorter than \(\varphi\), give \(\chi=\chi'\).
\item \((\exists v\chi)\) and \((\exists v'\chi')\) (or with \(\forall\)): the
  third symbols give \(v=v'\); \(\chi,\chi'\) occur at position \(4\), and
  (O) with the induction hypothesis gives \(\chi=\chi'\).
\item \((\chi\circ\theta)\) and \((\chi'\circ'\theta')\): \(\chi,\chi'\)
  occur at position \(2\), so \(\chi=\chi'\) by (O) and the induction
  hypothesis; the entry at position \(|\chi|+2\le|\varphi|\) gives
  \(\circ=\circ'\); \(\theta,\theta'\) occur at position \(|\chi|+3\), and
  again \(\theta=\theta'\).
\end{itemize}
In each case all components agree, so \(\varphi=\varphi'\).  Claims T and F
are clause~3.

\emph{Clause 1.}  Every member of \(P\) is a formula of \(L\), and the
formulas of \(L\) are closed under \(\neg\) and the binary connectives, so
by minimality every propositional formula of \(P\) is a formula of \(L\).
Conversely, let \(G\) be the set of formulas of \(L\) that are propositional
formulas of \(P\).  \(G\) contains every atomic formula, since atomic
formulas belong to \(P\).  If \(\varphi,\psi\in G\), then \(\neg\varphi\)
and \((\varphi\circ\psi)\) are propositional formulas of \(P\), so they lie
in \(G\).  If \(\varphi\in G\) and \(v\) is a variable, then
\((\exists v\varphi)\) and \((\forall v\varphi)\) are formulas of \(L\) that
belong to \(P\), so they lie in \(G\).  Thus \(G\) contains the atomic
formulas and is closed under the formation rules of \cref{def:p02-w02-p02-005-bw-syntax}, and by
minimality every formula of \(L\) lies in \(G\).

\emph{Clause 2.}  By inversion, a formula is a member of \(P\) (forms
\((t_1=t_2)\), \(R(t_1\cdots t_n)\), \((\exists v\varphi)\),
\((\forall v\varphi)\)), or of the form \(\neg\varphi\), or of the form
\((\varphi\circ\psi)\), and by the analysis of its first two symbols it is
of exactly one of these three kinds.  If \(\neg\varphi=\neg\varphi'\), then
\(\varphi=\varphi'\).  If \((\varphi\circ\psi)=(\varphi'\circ'\psi')\), then
\(\varphi,\varphi'\) occur at position \(2\) of the same expression, so by
(O) one is an initial segment of the other and Claim F gives
\(\varphi=\varphi'\); then \(\circ=\circ'\) is the entry at position
\(|\varphi|+2\), and \(\psi\) and \(\psi'\) both consist of the entries from
position \(|\varphi|+3\) to the second-last.
\end{proof}
\end{annexitem}

\begin{annexitem}{annex:p02-w02b-p02-001-recursion}{thm:p02-w02b-p02-001-recursion}{Recursion}
\begin{proof}
Let \(F\) be the set of propositional formulas of \(P\).  In each of the
three kinds the components are strictly shorter than the formula: \(B\) is
shorter than \(\neg B\), and \(B,D\) are shorter than \((B\circ D)\).

\emph{Uniqueness.}  Let \(h,h'\) both satisfy the three equations.  We show
\(h(A)=h'(A)\) by strong induction on \(|A|\).  If \(A\in P\), both values
are \(g_0(A)\).  If \(A=\neg B\), the induction hypothesis gives
\(h(B)=h'(B)\), so \(h(A)=g_\neg(h(B))=g_\neg(h'(B))=h'(A)\).  If
\(A=(B\circ D)\), likewise \(h(A)=g_\circ(h(B),h(D))=h'(A)\).

\emph{Existence.}  For \(n\geq0\) let \(F_n\) be the set of formulas of
length at most \(n\).  We define functions \(h_n\) on \(F_n\) by recursion
on \(n\), such that \(h_n\) satisfies the three equations for every
formula in \(F_n\) and \(h_{n+1}\) extends \(h_n\).  Put \(h_0=\emptyset\);
\(F_0\) is empty because formulas are nonempty.  Given \(h_n\), let
\(h_{n+1}\) agree with \(h_n\) on \(F_n\), and for \(A\) of length \(n+1\)
put \(h_{n+1}(A)=g_0(A)\) if \(A\in P\), \(h_{n+1}(A)=g_\neg(h_n(B))\) if
\(A=\neg B\), and \(h_{n+1}(A)=g_\circ(h_n(B),h_n(D))\) if
\(A=(B\circ D)\).  Because the formulas read uniquely, exactly one case
applies to \(A\) and its components \(B\), \(D\) and connective \(\circ\)
are determined by \(A\); they have length at most \(n\), so \(h_n\) is
defined at them.  Hence \(h_{n+1}\) is a well-defined function on
\(F_{n+1}\), and it satisfies the equations at every formula of length
\(n+1\), and at shorter formulas because it extends \(h_n\).  Let
\(h=\bigcup_n h_n\).  It is a function on \(F\), since every \(A\in F\) lies
in \(F_{|A|}\) and the \(h_n\) extend one another, and it satisfies the
three equations because each instance involves only formulas of length at
most \(|A|\), where \(h\) agrees with \(h_{|A|}\).

\emph{Barwise's clauses.}  Take \(Y=\{\mathbf t,\mathbf f\}\), \(g_0=v\),
\(g_\neg(\mathbf t)=\mathbf f\), \(g_\neg(\mathbf f)=\mathbf t\),
\(g_\land(x,y)=\mathbf t\) iff \(x=y=\mathbf t\),
\(g_\lor(x,y)=\mathbf t\) iff \(x=\mathbf t\) or \(y=\mathbf t\), and
\(g_\to(x,y)=\mathbf f\) iff \(x=\mathbf t\) and \(y=\mathbf f\).  The three
equations for these functions are exactly the clauses of
\cref{def:p02-w02-p02-005-bw-truth-assignment}, so \(\bar v\) exists and is
unique.
\end{proof}
\end{annexitem}

\begin{annexitem}{annex:p02-w02b-p02-001-coincidence}{thm:p02-w02b-p02-001-coincidence}{Coincidence}
\begin{proof}
Strong induction on \(|A|\), using the three kinds.  If \(A\in P\), its
only prime constituent is \(A\), so \(\bar v(A)=v(A)=w(A)=\bar w(A)\).  If
\(A=\neg B\), the prime constituents of \(B\) are those of \(A\), so
\(\bar v(B)=\bar w(B)\) by the induction hypothesis, and the clause for
\(\neg\) gives \(\bar v(A)=\bar w(A)\).  If \(A=(B\circ D)\), then \(v\) and
\(w\) agree on the prime constituents of \(B\) and on those of \(D\), since
both sets are contained in those of \(A\); the induction hypothesis gives
\(\bar v(B)=\bar w(B)\) and \(\bar v(D)=\bar w(D)\), and the clause for
\(\circ\) gives \(\bar v(A)=\bar w(A)\).
\end{proof}
\end{annexitem}

\begin{annexitem}{annex:p02-w02b-p02-001-tables}{thm:p02-w02b-p02-001-truth-table}{Truth tables}
\begin{proof}
Write \(Q=\{q_1,\ldots,q_k\}\).  By induction on length every formula has
at least one prime constituent (a member of \(P\) has itself, and the
other two kinds inherit those of their components), so \(k\geq1\).

\emph{Every row has an extension.}  Given a row \(r\), let \(v(q)=r(q)\)
for \(q\in Q\) and \(v(p)=\mathbf f\) for \(p\in P-Q\).  This is a truth
assignment extending \(r\).

\emph{The value at a row is well defined.}  Two truth assignments extending
\(r\) agree on \(Q\), the set of prime constituents of \(A\), so by
\cref{thm:p02-w02b-p02-001-coincidence} they give \(A\) the same value.

\emph{Number of rows.}  A row chooses one of two values at each of the
\(k\) distinct points of \(Q\), independently, so there are \(2^k\) rows.

\emph{Tautology.}  If \(A\) is a tautology, then \(\bar v(A)=\mathbf t\) for
every truth assignment \(v\), in particular for an extension of each row,
so every row has value \(\mathbf t\).  Conversely, suppose every row has
value \(\mathbf t\), and let \(v\) be any truth assignment.  Its restriction
\(r\) to \(Q\) is a row and \(v\) extends \(r\), so
\(\bar v(A)\) is the value at \(r\), which is \(\mathbf t\).  Hence \(A\) is
a tautology.

\emph{Consistency.}  If \(\bar v(A)=\mathbf t\) for some \(v\), the row
\(v\) restricted to \(Q\) has value \(\mathbf t\).  If some row \(r\) has
value \(\mathbf t\), its extension constructed above is a truth assignment
\(v\) with \(\bar v(A)=\mathbf t\), so \(A\) is consistent.
\end{proof}
\end{annexitem}

\section{Proofs for the completeness of the cut-free system G}
\label{sec:annex-p02-w02b-p02-001-g-completeness}

The notation is that of \cref{sec:p02-w02b-p02-001-g-completeness}; rule
instances are read as in \cref{def:p02-w02b-p02-001-g-reading}.  For a
formula \(\chi\), \(\#\chi\) is the number of occurrences of connectives and
quantifiers in \(\chi\).

\begin{annexitem}{annex:p02-w02b-p02-001-g-substitution}{thm:p02-w02b-p02-001-g-substitution}{Clean substitution}
\begin{proof}
\emph{Terms.}  For a term \(\tau\), \(\tau(t/v)\) replaces each occurrence
of \(v\) in \(\tau\) by \(t\).  By induction on \(\tau\):
\(\operatorname{Var}(\tau(t/v))\) is \(\operatorname{Var}(\tau)\) if
\(v\notin\operatorname{Var}(\tau)\) and
\((\operatorname{Var}(\tau)-\{v\})\cup\operatorname{Var}(t)\) otherwise
(the cases \(\tau=v\), \(\tau\) a variable other than \(v\), \(\tau\) a
constant symbol are immediate, and for \(f(\tau_1\cdots\tau_n)\) take the
union over the arguments); and \(\tau(t/v)=\tau\) if
\(v\notin\operatorname{Var}(\tau)\).

\emph{Part 1.}  Induction on \(\theta\), following the clauses of
\cref{def:p02-w02b-p02-001-g-reading}(2).  For atomic \(\theta\) there are
no quantifiers or connectives, and the free variables are the variables of
the terms, so the term case applies.  For \(\neg\sigma\) and
\((\sigma_1\circ\sigma_2)\) the operation is applied to the components,
which satisfy the hypothesis; the claims about quantifier occurrences,
scopes and the count follow from those for the components, and the free
variables are the union, so if \(v\) is free in some components and not in
others the formula
\((\FV(\theta)-\{v\})\cup\operatorname{Var}(t)\) results.  For
\((Qv\sigma)\) the formula is unchanged and \(v\notin\FV(Qv\sigma)\), as
claimed.  For \((Qw\sigma)\) with \(w\neq v\) the result is
\((Qw\,\sigma(t/v))\): the outer quantifier keeps its variable, its scope is
the image of \(\sigma\), and the induction hypothesis applies to \(\sigma\).
Since \(w\) occurs bound in \(\theta\), \(w\notin\operatorname{Var}(t)\).
So \(\FV(Qw\,\sigma(t/v))=\FV(\sigma(t/v))-\{w\}\), which is
\(\FV(\sigma)-\{w\}=\FV(Qw\sigma)\) when \(v\notin\FV(\sigma)\), and
\(((\FV(\sigma)-\{v\})\cup\operatorname{Var}(t))-\{w\}
=(\FV(Qw\sigma)-\{v\})\cup\operatorname{Var}(t)\) otherwise; here
\(v\in\FV(\sigma)\) iff \(v\in\FV(Qw\sigma)\) because \(w\neq v\).

\emph{Part 2.}  The substitution \((u/y)\) is clean because \(u\) does not
occur in \(\theta\).  By part 1 the quantified variables of \(\theta(u/y)\)
are those of \(\theta\), so \(y\) is not bound in \(\theta(u/y)\) and
\((y/u)\) is clean.  For a term \(\tau\) without \(u\),
\(\tau(u/y)(y/u)=\tau\) by induction: \(y\) goes to \(u\) and back to
\(y\); a variable \(w\neq y\) is also \(\neq u\) and is fixed by both steps;
constants are fixed; function terms follow from their arguments.  For
formulas use induction on \(\theta\); atomic formulas and connectives
reduce to their parts.  For \((Qw\sigma)\) in \(\theta\), \(w\neq y\)
because \(y\) is not bound in \(\theta\), and \(w\neq u\) because \(u\)
does not occur in \(\theta\); so \((Qw\sigma)(u/y)=(Qw\,\sigma(u/y))\) and
\((Qw\,\sigma(u/y))(y/u)=(Qw\,\sigma(u/y)(y/u))=(Qw\sigma)\) by the
induction hypothesis for \(\sigma\).

\emph{Part 3.}  Induction on \(\tau\).  If \(\tau=x\), both sides are
\(w(r/z)\).  If \(\tau\) is a variable \(y\neq x\), then \(y\neq z\)
because \(z\) does not occur in \(\tau\), so both sides are \(y\).  A
constant symbol is fixed on both sides, and \(f(\tau_1\cdots\tau_n)\)
follows from its arguments, none of which contains \(z\).  (If \(x=z\) the
first case cannot arise, and both sides equal \(\tau\).)  For an atomic
formula, substitution acts on each of its terms, so the term case applies
to each.

\emph{Part 4.}  First, for a term \(\tau\),
\(\tau(t/v)^{\mathfrak M}(s)=\tau^{\mathfrak M}(s(\tfrac bv))\), by
induction: for \(\tau=v\) both sides are \(b\); for a variable \(w\neq v\)
both are \(s(w)\); constants are immediate; function terms follow by
applying \(f^{\mathfrak M}\).  Second, if assignments agree on
\(\operatorname{Var}(\tau)\), they give \(\tau\) the same value, again by
induction on \(\tau\); this is the fact Barwise notes after the definition
of \(t^{\mathfrak M}\) \cite[ch.~A.1, \S3, p.~21]{Barwise1977FirstOrder}.

Now induct on \(\theta\), for all assignments \(s\) at once.  Atomic
\(\theta\): by the term fact, the values of the terms of \(\theta(t/v)\)
at \(s\) are those of the terms of \(\theta\) at \(s(\tfrac bv)\), and
clauses (i), (ii) of \cref{def:p02-w02-p02-005-bw-satisfaction} decide
both sides the same way.  Connectives: the satisfaction clauses and the
induction hypothesis.  For \((Qv\sigma)\): the formula is unchanged, and
\(s\), \(s(\tfrac bv)\) agree on \(\FV(Qv\sigma)\), which does not contain
\(v\); so \cref{lem:fol-coincidence} gives the claim.  For \((Qw\sigma)\)
with \(w\neq v\): \(w\notin\operatorname{Var}(t)\), as \(w\) is bound in
\(\theta\).  For \(a\in M\), the assignments \(s(\tfrac aw)\) and \(s\)
agree on \(\operatorname{Var}(t)\), so \(t^{\mathfrak M}(s(\tfrac aw))=b\);
the induction hypothesis for \(\sigma\) at \(s(\tfrac aw)\) gives
\(\mathfrak M\models\sigma(t/v)[s(\tfrac aw)]\) iff
\(\mathfrak M\models\sigma[s(\tfrac aw)(\tfrac bv)]\), and
\(s(\tfrac aw)(\tfrac bv)=s(\tfrac bv)(\tfrac aw)\) since \(w\neq v\).
Hence \(\mathfrak M\models(\forall w\,\sigma(t/v))[s]\) iff
\(\mathfrak M\models\sigma[s(\tfrac bv)(\tfrac aw)]\) for every \(a\in M\),
iff \(\mathfrak M\models(\forall w\sigma)[s(\tfrac bv)]\); for \(\exists\)
replace ``every'' by ``some''.
\end{proof}
\end{annexitem}

\begin{annexitem}{annex:p02-w02b-p02-001-g-completeness}{thm:p02-w02b-p02-001-g-completeness}{Completeness of the cut-free system G}
\begin{proof}
\emph{Soundness.}  A derivation in \(\mathbf G\) is a derivation in
\(\mathbf G^{+}\).  The proof of \cref{thm:p02-w02-p02-005-bw-completeness-g}
shows that every axiom of \(\mathbf G^{+}\) is valid and that every rule of
\(\mathbf G^{+}\) passes from valid premises to a valid conclusion; by
induction on the height of a derivation tree its end sequent is valid.

\emph{Completeness.}  Let \(\Gamma_0\vdash\Delta_0\) separate its variables
and have no derivation in \(\mathbf G\); we construct a structure and an
assignment satisfying every member of \(\Gamma_0\) and no member of
\(\Delta_0\), so that \(\Gamma_0\vdash\Delta_0\) is not valid.  ``Derivable''
means ``has a derivation in \(\mathbf G\)''.  Two facts about derivations
are used throughout.  (F1) An axiom is derivable: the one-node tree is a
derivation.  (F2) If every premise of an instance of a rule of \(\mathbf G\)
is derivable, so is its conclusion: place the derivations of the premises
above it.

\emph{Step 0: clean formulas.}  Let \(B\), \(L_0\) and \(U\) be as in
\cref{def:p02-w02b-p02-001-g-search}; \(B\) is finite, and \(U\) contains
every variable outside \(B\).  Call a formula \emph{\(B\)-clean} if it is a
formula of \(L_0\), every variable quantified in it lies in \(B\), no
variable of \(B\) occurs free in it, and no quantifier on a variable occurs
in the scope of a quantifier on the same variable.  Every member of
\(\Gamma_0\cup\Delta_0\) is \(B\)-clean by hypothesis.  Let \(\chi\) be
\(B\)-clean.
\begin{enumerate}
\item[(i)] The components of \(\neg\varphi\) and \((\varphi\circ\psi)\) are
      \(B\)-clean: their quantifier occurrences and scopes are among those
      of \(\chi\), and their free variables are free in \(\chi\).
\item[(ii)] If \(\chi\) is \((\forall v\theta)\) or \((\exists v\theta)\)
      and \(t\in U\), then \(\theta(t/v)\) is clean, since the variables
      bound in \(\theta\) lie in \(B\) and those of \(t\) do not; and
      \(\theta(t/v)\) is \(B\)-clean.  Indeed, by part~1 of
      \cref{thm:p02-w02b-p02-001-g-substitution} its quantifier occurrences
      and scopes are those of \(\theta\), so the first and third conditions
      hold; its free variables lie in
      \((\FV(\theta)-\{v\})\cup\operatorname{Var}(t)\), and
      \(\FV(\theta)-\{v\}\subseteq\FV(\chi)\) is disjoint from \(B\), as is
      \(\operatorname{Var}(t)\).  It is a formula of \(L_0\) because \(t\)
      is a term of \(L_0\).
\item[(iii)] In \(\chi=(Qy\theta)\) the variable \(y\) has no bound
      occurrence in \(\theta\) (no quantifier on \(y\) lies in the scope of
      the outer one), so every occurrence of \(y\) in \(\theta\) is free.
\item[(iv)] An atomic formula of \(L_0\) whose terms lie in \(U\) is
      \(B\)-clean; conversely the terms of a \(B\)-clean atomic formula lie
      in \(U\), since all its variables are free.
\end{enumerate}

\emph{Step 1: tasks and schedule.}  Variables are taken in Barwise's
infinite sequence \(x,y,z,x_1,y_1,\ldots\)
(\cref{def:p02-w02-p02-005-bw-language}); ``the first variable'' with a
property refers to this order.  Formulas of \(L_0\), terms in \(U\) and
variables are finite sequences of symbols from a countable set, namely
\(L_0\), the logical symbols and the variables; the set of finite
sequences over a countable set is countable.  So the set of tasks of
\cref{def:p02-w02b-p02-001-g-search} is countable, and it is infinite.  Fix
a list \(\tau_0,\tau_1,\ldots\) in which every task occurs.  For
\(n\geq0\) write \(n+1=2^{k}(2m+1)\) with \(k,m\geq0\), and set
\(\pi(n)=k\).  For each \(k\), \(\pi(n)=k\) for the infinitely many
\(n=2^{k}(2m+1)-1\).  Stage \(n\) works on the task \(\tau_{\pi(n)}\).

\emph{Step 2: the stages.}  We define sequents \(\Gamma_n\vdash\Delta_n\)
with the invariants (I1) \(\Gamma_n,\Delta_n\) are finite sets of
\(B\)-clean formulas, (I2) \(\Gamma_n\subseteq\Gamma_{n+1}\) and
\(\Delta_n\subseteq\Delta_{n+1}\), (I3) \(\Gamma_n\vdash\Delta_n\) is not
derivable.  Stage \(0\) is the given sequent.  Suppose
\(\Gamma_n\vdash\Delta_n\) satisfies (I1) and (I3).  It has no formula on
both sides, since otherwise it is an instance \(\Gamma,\chi\vdash\Delta,\chi\)
of the axiom and is derivable by (F1).  In each case below, the new
sequent is the premise, or one of the premises, of a rule instance whose
conclusion is \(\Gamma_n\vdash\Delta_n\) itself, with side sets
\(\Gamma=\Gamma_n\) and \(\Delta=\Delta_n\); this is possible because the
principal formula already lies in \(\Gamma_n\) or \(\Delta_n\)
(\cref{def:p02-w02b-p02-001-g-reading}(1)).  By (F2) and (I3) for stage
\(n\), a one-premise instance has an underivable premise, and a two-premise
instance has at least one underivable premise; we take the first premise
listed if it is underivable and the second otherwise.  So (I3) holds at
stage \(n+1\), and (I2) holds because every premise contains
\(\Gamma_n\) and \(\Delta_n\).

\emph{Task \((\chi,t)\) with \(\chi\in\Gamma_n\).}
\begin{itemize}
\item \(\chi\) atomic: \(\Gamma_{n+1}=\Gamma_n\), \(\Delta_{n+1}=\Delta_n\).
\item \(\chi=\neg\varphi\): rule \((\neg\vdash)\); the premise is
      \(\Gamma_n\vdash\Delta_n,\varphi\).
\item \(\chi=(\varphi\land\psi)\): rule \((\land\vdash)\); the premise is
      \(\Gamma_n,\varphi,\psi\vdash\Delta_n\).
\item \(\chi=(\varphi\lor\psi)\): rule \((\lor\vdash)\); the premises are
      \(\Gamma_n,\varphi\vdash\Delta_n\) and \(\Gamma_n,\psi\vdash\Delta_n\).
\item \(\chi=(\varphi\to\psi)\): rule \((\to\vdash)\); the premises are
      \(\Gamma_n\vdash\Delta_n,\varphi\) and \(\Gamma_n,\psi\vdash\Delta_n\).
\item \(\chi=(\forall v\theta)\): rule \((\forall\vdash)\) with
      \(\varphi(v)=\theta\) and the term \(t\); the premise is
      \(\Gamma_n,\theta(t/v)\vdash\Delta_n\), and the substitution is clean
      by Step~0(ii).
\item \(\chi=(\exists y\theta)\): let \(u\) be the first variable outside
      \(B\) that occurs in no member of \(\Gamma_n\cup\Delta_n\) (it exists
      because these sets are finite and infinitely many variables lie
      outside \(B\)).  Use rule \((\exists\vdash)\) with \(v=u\) and
      \(\varphi(u)=\theta(u/y)\); the premise is
      \(\Gamma_n,\theta(u/y)\vdash\Delta_n\).  This is an instance: \(u\)
      does not occur free in \(\Gamma_n\cup\Delta_n\); \(u\) does not occur
      in \(\theta\) and \(y\) is not bound in \(\theta\) (Step~0(iii)), so
      \(\varphi(u)(y/u)=\theta\) by part~2 of
      \cref{thm:p02-w02b-p02-001-g-substitution}, and the conclusion
      formula \(\exists y\,\varphi(y)\) is \(\chi\); finally \(y\neq u\)
      (as \(y\in B\)) and \(y\) does not occur in \(\varphi(u)\), because
      every occurrence of \(y\) in \(\theta\) is free (Step~0(iii)) and was
      replaced by \(u\), so the reading
      \cref{def:p02-w02b-p02-001-g-reading}(4) is met.
\end{itemize}

\emph{Task \((\chi,t)\) with \(\chi\in\Delta_n\).}  The same with the right
rules: \((\vdash\neg)\) gives \(\Gamma_n,\varphi\vdash\Delta_n\);
\((\vdash\land)\) gives the premises \(\Gamma_n\vdash\Delta_n,\varphi\) and
\(\Gamma_n\vdash\Delta_n,\psi\); \((\vdash\lor)\) gives
\(\Gamma_n\vdash\Delta_n,\varphi,\psi\); \((\vdash\to)\) gives
\(\Gamma_n,\varphi\vdash\Delta_n,\psi\); \((\vdash\exists)\) with
\(\varphi(v)=\theta\) and the term \(t\) gives
\(\Gamma_n\vdash\Delta_n,\theta(t/v)\); and for \(\chi=(\forall y\theta)\),
\((\vdash\forall)\) with \(u\) chosen as above gives
\(\Gamma_n\vdash\Delta_n,\theta(u/y)\), the verification being the one
just given.  An atomic \(\chi\) changes nothing.

\emph{Task \((\chi,t)\) with \(\chi\notin\Gamma_n\cup\Delta_n\).}  Nothing
changes.

\emph{Task \((a,b,\psi,x)\).}  If \(\Gamma_n\) contains
\(E=(a=b)\) or \(E=(b=a)\) (if both, take \((a=b)\)) and
\(\psi(b/x)\in\Delta_n\), let \(z\) be the first variable outside \(B\)
that occurs in none of \(\psi\), \(a\), \(b\), and put
\(\psi^{*}(z)=\psi(z/x)\) and \(\varphi^{*}(z)=E\).  By part~3 of
\cref{thm:p02-w02b-p02-001-g-substitution}, \(\psi^{*}(r/z)=\psi(r/x)\) for
\(r=a,b\); since \(z\) does not occur in \(E\),
\(\varphi^{*}(a/z)=\varphi^{*}(b/z)=E\).  The equality rule with
\(t_1=a\), \(t_2=b\), this \(E\), \(\Gamma=\Gamma_n\), \(\Delta=\Delta_n\)
has premise \(\Gamma_n,E\vdash\Delta_n,\psi(a/x)\), that is
\(\Gamma_n\vdash\Delta_n,\psi(a/x)\), and conclusion
\(\Gamma_n,E,E\vdash\Delta_n,\psi(b/x)\), that is
\(\Gamma_n\vdash\Delta_n\); the new sequent is this premise.  The
substitutions are clean, as atomic formulas have no bound variables.
Otherwise nothing changes.

Every new formula is \(B\)-clean: components and instances by Step~0(i)
and~(ii) (in the \(\exists\) and \(\forall\) cases \(u\in U\)), and
\(\psi(a/x)\) by Step~0(iv), since its terms are obtained from terms in
\(U\) by replacing \(x\) with \(a\in U\).  So (I1) holds at stage \(n+1\).

\emph{Step 3: the limit.}  Let \(\Gamma_\omega=\bigcup_n\Gamma_n\) and
\(\Delta_\omega=\bigcup_n\Delta_n\).
\begin{enumerate}
\item[(L1)] \(\Gamma_\omega\cap\Delta_\omega=\emptyset\).  If
      \(\chi\in\Gamma_m\) and \(\chi\in\Delta_{m'}\), then by (I2)
      \(\chi\in\Gamma_n\cap\Delta_n\) for \(n=\max(m,m')\), which was
      excluded in Step~2.
\item[(L2)] No formula \((r=r)\) lies in \(\Delta_\omega\): otherwise some
      \(\Delta_n\) contains it, and \(\Gamma_n\vdash\Delta_n\) is an
      instance \(\Gamma\vdash\Delta,(r=r)\) of the equality axiom, derivable
      by (F1).
\item[(L3)] If \(\chi\in\Gamma_\omega\), then: for \(\chi=\neg\varphi\),
      \(\varphi\in\Delta_\omega\); for \((\varphi\land\psi)\), both
      components lie in \(\Gamma_\omega\); for \((\varphi\lor\psi)\), one of
      them does; for \((\varphi\to\psi)\), \(\varphi\in\Delta_\omega\) or
      \(\psi\in\Gamma_\omega\); for \((\forall v\theta)\),
      \(\theta(t/v)\in\Gamma_\omega\) for every \(t\in U\); for
      \((\exists y\theta)\), \(\theta(u/y)\in\Gamma_\omega\) for some
      variable \(u\notin B\).  Dually, if \(\chi\in\Delta_\omega\): for
      \(\neg\varphi\), \(\varphi\in\Gamma_\omega\); for
      \((\varphi\land\psi)\), one component lies in \(\Delta_\omega\); for
      \((\varphi\lor\psi)\), both do; for \((\varphi\to\psi)\),
      \(\varphi\in\Gamma_\omega\) and \(\psi\in\Delta_\omega\); for
      \((\exists v\theta)\), \(\theta(t/v)\in\Delta_\omega\) for every
      \(t\in U\); for \((\forall y\theta)\), \(\theta(u/y)\in\Delta_\omega\)
      for some variable \(u\notin B\).
\item[(L4)] If \((a=b)\) or \((b=a)\) lies in \(\Gamma_\omega\) and
      \(\psi(b/x)\in\Delta_\omega\), where \((a,b,\psi,x)\) is a task, then
      \(\psi(a/x)\in\Delta_\omega\).
\end{enumerate}
For (L3), let \(\chi\in\Gamma_m\), say, fix \(t\in U\), and let \(k\) be
an index with \(\tau_k=(\chi,t)\).  Choose \(n\geq m\) with \(\pi(n)=k\).
By (I2), \(\chi\in\Gamma_n\), so stage \(n\) applies the case for
\(\chi\) and puts the listed formula, or one of the two listed formulas,
into \(\Gamma_{n+1}\) or \(\Delta_{n+1}\).  For \(\forall\) on the left and
\(\exists\) on the right this is done for each \(t\in U\).  For (L4), take
\(m\) with the equation in \(\Gamma_m\) and \(\psi(b/x)\in\Delta_m\), let
\(\tau_k=(a,b,\psi,x)\), and choose \(n\geq m\) with \(\pi(n)=k\); the
conditions of the task hold at stage \(n\), so
\(\psi(a/x)\in\Delta_{n+1}\).

\emph{Step 4: interchangeable terms.}  For \(a,b\in U\) let
\(a\approx b\) mean: for every atomic formula \(\psi\) of \(L_0\) with terms
in \(U\) and every variable \(x\notin B\),
\(\psi(a/x)\in\Delta_\omega\) iff \(\psi(b/x)\in\Delta_\omega\).
\begin{enumerate}
\item[(i)] \(\approx\) is an equivalence relation on \(U\), since it is
      defined by a family of biconditionals.
\item[(ii)] If \((a=b)\in\Gamma_\omega\), then \(a\approx b\).  Given
      \(\psi,x\): if \(\psi(b/x)\in\Delta_\omega\), (L4) for the task
      \((a,b,\psi,x)\) gives \(\psi(a/x)\in\Delta_\omega\); if
      \(\psi(a/x)\in\Delta_\omega\), (L4) for the task \((b,a,\psi,x)\),
      whose condition ``\((b=a)\) or \((a=b)\) lies in \(\Gamma_\omega\)''
      holds, gives \(\psi(b/x)\in\Delta_\omega\).
\item[(iii)] \(\approx\) is a congruence: if \(f\in L_0\) is \(n\)-ary,
      \(r_1,\ldots,r_n\in U\) and \(a\approx b\), then
      \(f(r_1\cdots r_{i-1}\,a\,r_{i+1}\cdots r_n)\approx
      f(r_1\cdots r_{i-1}\,b\,r_{i+1}\cdots r_n)\).  Given \(\psi,x\), let
      \(z\) be a variable outside \(B\) occurring in none of \(\psi\),
      \(a\), \(b\), \(r_1,\ldots,r_n\), let
      \(w=f(r_1\cdots r_{i-1}\,z\,r_{i+1}\cdots r_n)\), and let
      \(\psi''=\psi(w/x)\), an atomic formula of \(L_0\) with terms in
      \(U\).  By part~3 of \cref{thm:p02-w02b-p02-001-g-substitution},
      \(\psi''(a/z)=\psi(w(a/z)/x)\), and
      \(w(a/z)=f(r_1\cdots a\cdots r_n)\) because \(z\) occurs in \(w\) only
      at position \(i\); likewise for \(b\).  So
      \(\psi(f(\cdots a\cdots)/x)\in\Delta_\omega\) iff
      \(\psi''(a/z)\in\Delta_\omega\) iff \(\psi''(b/z)\in\Delta_\omega\)
      (as \(a\approx b\) and \(z\notin B\)) iff
      \(\psi(f(\cdots b\cdots)/x)\in\Delta_\omega\).  Changing the
      arguments one at a time and using transitivity, \(a_j\approx b_j\)
      for all \(j\) gives
      \(f(a_1\cdots a_n)\approx f(b_1\cdots b_n)\).
\item[(iv)] If \((a=b)\in\Delta_\omega\), then not \(a\approx b\).  Let
      \(x\notin B\) occur in neither \(a\) nor \(b\) and let
      \(\psi=(a=x)\).  Then \(\psi(b/x)=(a=b)\in\Delta_\omega\); if
      \(a\approx b\), then \(\psi(a/x)=(a=a)\in\Delta_\omega\), against
      (L2).
\item[(v)] If \(R(a_1\cdots a_n)\in\Delta_\omega\) and
      \(R(a_1'\cdots a_n')\in\Gamma_\omega\), then \(a_j\approx a_j'\)
      fails for some \(j\).  Suppose \(a_j\approx a_j'\) for all \(j\).  We
      show by induction on \(i\leq n\) that
      \(R(a_1'\cdots a_i'\,a_{i+1}\cdots a_n)\in\Delta_\omega\).  For
      \(i=0\) this is the hypothesis.  Given it for \(i<n\), let
      \(x\notin B\) occur in none of the \(a_j,a_j'\), and let
      \(\psi=R(a_1'\cdots a_i'\,x\,a_{i+2}\cdots a_n)\).  Then
      \(\psi(a_{i+1}/x)\in\Delta_\omega\), so
      \(\psi(a_{i+1}'/x)\in\Delta_\omega\), which is the claim for
      \(i+1\).  For \(i=n\), \(R(a_1'\cdots a_n')\) lies in
      \(\Delta_\omega\cap\Gamma_\omega\), against (L1).  (The terms here lie
      in \(U\) by Step~0(iv).)
\end{enumerate}

\emph{Step 5: the structure.}  Write \([a]\) for the class of \(a\in U\)
under \(\approx\), and let \(M=\{[a]:a\in U\}\), a nonempty set.  Let
\(x^{\circ}\) be the first variable outside \(B\).  Define a structure
\(\mathfrak M\) for \(L\) with domain \(M\): for a constant symbol
\(c\in L_0\), \(c^{\mathfrak M}=[c]\); for an \(n\)-ary function symbol
\(f\in L_0\), \(f^{\mathfrak M}([a_1],\ldots,[a_n])=[f(a_1\cdots a_n)]\),
which is well defined by Step~4(iii); for an \(n\)-ary relation symbol
\(R\in L_0\), \(R^{\mathfrak M}=\{([a_1],\ldots,[a_n]):
R(a_1\cdots a_n)\in\Gamma_\omega\}\).  Symbols of \(L\) outside \(L_0\)
are interpreted arbitrarily, say constants by \([x^{\circ}]\), functions
by the constant function with value \([x^{\circ}]\), relations by
\(\emptyset\).  Let \(s(w)=[w]\) for a variable \(w\notin B\) and
\(s(w)=[x^{\circ}]\) for \(w\in B\).  By induction on terms,
\(a^{\mathfrak M}(s)=[a]\) for every \(a\in U\); so every element of \(M\)
is the value at \(s\) of a term in \(U\).

For atomic members: if \((a=b)\in\Gamma_\omega\), then \([a]=[b]\) by
Step~4(ii), so \(s\) satisfies \((a=b)\); if
\(R(a_1\cdots a_n)\in\Gamma_\omega\), then
\(([a_1],\ldots,[a_n])\in R^{\mathfrak M}\), so \(s\) satisfies it.  If
\((a=b)\in\Delta_\omega\), then \([a]\neq[b]\) by Step~4(iv), so \(s\) does
not satisfy it.  If \(R(a_1\cdots a_n)\in\Delta_\omega\) and \(s\)
satisfied it, then \(([a_1],\ldots,[a_n])\in R^{\mathfrak M}\), so some
\(R(a_1'\cdots a_n')\in\Gamma_\omega\) has \(a_j'\approx a_j\) for every
\(j\), against Step~4(v).

\emph{Step 6: truth lemma.}  For every \(\chi\in\Gamma_\omega\cup\Delta_\omega\):
if \(\chi\in\Gamma_\omega\) then \(\mathfrak M\models\chi[s]\), and if
\(\chi\in\Delta_\omega\) then not \(\mathfrak M\models\chi[s]\).  The proof
is by induction on \(\#\chi\).  The atomic case is Step~5.  The
propositional cases follow from (L3), the induction hypothesis for the
components, and the satisfaction clauses; for instance, if
\((\varphi\to\psi)\in\Gamma_\omega\), then \(\varphi\in\Delta_\omega\), so
\(s\) does not satisfy \(\varphi\), or \(\psi\in\Gamma_\omega\), so \(s\)
satisfies \(\psi\); either way \(s\) satisfies \((\varphi\to\psi)\).  For
quantifiers note that \(\#\theta(t/v)=\#\theta<\#\chi\) by part~1 of
\cref{thm:p02-w02b-p02-001-g-substitution}.
\begin{itemize}
\item \((\forall v\theta)\in\Gamma_\omega\): for every \(t\in U\),
      \(\theta(t/v)\in\Gamma_\omega\), so \(\mathfrak M\models\theta(t/v)[s]\),
      so by part~4 of \cref{thm:p02-w02b-p02-001-g-substitution}
      \(\mathfrak M\models\theta[s(\tfrac{[t]}{v})]\).  Every element of
      \(M\) is some \([t]\), so \(\mathfrak M\models(\forall v\theta)[s]\).
\item \((\exists y\theta)\in\Gamma_\omega\): some
      \(\theta(u/y)\in\Gamma_\omega\), so
      \(\mathfrak M\models\theta[s(\tfrac{[u]}{y})]\) by part~4, and
      \(\mathfrak M\models(\exists y\theta)[s]\).
\item \((\exists v\theta)\in\Delta_\omega\): for every \(t\in U\),
      \(\theta(t/v)\in\Delta_\omega\), so \(s(\tfrac{[t]}{v})\) does not
      satisfy \(\theta\); as every element is some \([t]\),
      \(s\) does not satisfy \((\exists v\theta)\).
\item \((\forall y\theta)\in\Delta_\omega\): some
      \(\theta(u/y)\in\Delta_\omega\), so \(s(\tfrac{[u]}{y})\) does not
      satisfy \(\theta\), and \(s\) does not satisfy \((\forall y\theta)\).
\end{itemize}

\emph{Step 7: conclusion.}  Since \(\Gamma_0\subseteq\Gamma_\omega\) and
\(\Delta_0\subseteq\Delta_\omega\), the assignment \(s\) in the structure
\(\mathfrak M\) for \(L\) satisfies every member of \(\Gamma_0\) and no
member of \(\Delta_0\).  So \(\Gamma_0\vdash\Delta_0\) is not valid.  By
contraposition, every valid sequent that separates its variables is
derivable in \(\mathbf G\).
\end{proof}
\end{annexitem}

\begin{annexitem}{annex:p02-w02b-p02-001-g-cutfree}{thm:p02-w02b-p02-001-g-cut-elimination}{Cut removal}
\begin{proof}
\emph{Part 1.}  Every rule of \(\mathbf G\) is a rule of \(\mathbf G^{+}\),
so a derivation in \(\mathbf G\) is one in \(\mathbf G^{+}\).  Conversely,
let a sequent that separates its variables have a derivation in
\(\mathbf G^{+}\).  By the soundness of \(\mathbf G^{+}\) (first paragraph
of the proof of \cref{thm:p02-w02b-p02-001-g-completeness}) it is valid,
and by \cref{thm:p02-w02b-p02-001-g-completeness} it has a derivation in
\(\mathbf G\).

\emph{Part 2.}  Suppose \(T\models\varphi\).  By
\cref{thm:p02-w02-p02-005-bw-completeness-g}, \(\varphi\) is derivable from
\(T\) in \(\mathbf G^{+}\): some sequent \(\Gamma\vdash\{\varphi\}\) with
finite \(\Gamma\subseteq T\) has a derivation in \(\mathbf G^{+}\).  This
sequent consists of sentences, so it satisfies condition (a) of
\cref{def:p02-w02b-p02-001-g-valid}, and it satisfies condition (b) by the
hypothesis on \(T\) and \(\varphi\).  By part~1 it has a derivation in
\(\mathbf G\).  Conversely, a derivation in \(\mathbf G\) of
\(\Gamma\vdash\{\varphi\}\) with finite \(\Gamma\subseteq T\) is a
derivation in \(\mathbf G^{+}\), so \(\varphi\) is derivable from \(T\) in
\(\mathbf G^{+}\), and \(T\models\varphi\) by
\cref{thm:p02-w02-p02-005-bw-completeness-g}.
\end{proof}
\end{annexitem}

\subsection{Sequents that do not separate their variables}
\label{sec:annex-p02-w02b-p02-001-g-variants}

The notation is that of \cref{sec:p02-w02b-p02-001-g-variants}.  For a
formula \(\chi\), \(\operatorname{QV}(\chi)\) is the set of variables
quantified in \(\chi\).  By \cref{def:p02-w02-p02-005-bw-syntax},
\(\FV(\neg\alpha)=\FV(\alpha)\),
\(\FV(\alpha\circ\beta)=\FV(\alpha)\cup\FV(\beta)\) and
\(\FV(Qv\alpha)=\FV(\alpha)-\{v\}\); likewise
\(\operatorname{QV}(\neg\alpha)=\operatorname{QV}(\alpha)\),
\(\operatorname{QV}(\alpha\circ\beta)=\operatorname{QV}(\alpha)\cup\operatorname{QV}(\beta)\)
and \(\operatorname{QV}(Qv\alpha)=\operatorname{QV}(\alpha)\cup\{v\}\).
\emph{Occurrences of subformulas} are taken recursively: the occurrences in
\(\chi\) are \(\chi\) itself and the occurrences in the immediate
constituents of \(\chi\), namely \(\alpha\) for \(\neg\alpha\) and for
\((Qv\alpha)\), and \(\alpha,\beta\) for \((\alpha\circ\beta)\).  These
constituents are uniquely determined: for \(\neg\alpha\) and
\((\alpha\circ\beta)\) by clause~2 of
\cref{thm:p02-w02b-p02-001-fo-skeleton}, and for \((Qv\alpha)\) because
\(Q\) and \(v\) are its second and third symbols and \(\alpha\) is what
lies between them and the final right parenthesis.  Replacing one
occurrence inside an immediate constituent of \(\chi\) replaces it in that
constituent and leaves the rest of \(\chi\) as it was.  Hence:

\emph{(V)} If \(\alpha'\) is a variant of \(\alpha\), then \(\neg\alpha'\),
\((\alpha'\circ\beta)\), \((\beta\circ\alpha')\) and \((Qv\alpha')\) are
variants of \(\neg\alpha\), \((\alpha\circ\beta)\), \((\beta\circ\alpha)\) and
\((Qv\alpha)\) respectively: each renaming step inside \(\alpha\), forwards
or backwards, is a renaming step at the corresponding occurrence inside the
compound formula.

\begin{annexitem}{annex:p02-w02b-p02-001-g-clean-failure}{thm:p02-w02b-p02-001-g-clean-failure}{Cut cannot be eliminated under the reading without renaming}
\begin{proof}
Write \(\chi_0=\forall x\,\forall y\,P(xy)\) and
\(\chi_1=\forall x\,\forall z\,P(xz)\).  Rule instances are read as in
\cref{def:p02-w02b-p02-001-g-reading}.

\emph{The conditions on variables.}  The variable \(y\) is quantified in
\(\chi_0\) and free in \(P(yy)\), so condition (a) of
\cref{def:p02-w02b-p02-001-g-valid} fails.  The only quantifier occurrences
in the sequent are \(\forall x\) and \(\forall y\) in \(\chi_0\), on
different variables, so condition (b) holds.

\emph{Validity.}  Let \(\mathfrak M\models\chi_0[s]\) and put
\(a=s(y)\).  By the clause for \(\forall\) of
\cref{def:p02-w02-p02-005-bw-satisfaction} applied twice,
\(\mathfrak M\models P(xy)[s(\tfrac ax)(\tfrac ay)]\), so
\((a,a)\in P^{\mathfrak M}\), that is, \(\mathfrak M\models P(yy)[s]\).

\emph{A derivation in \(\mathbf G^{+}\).}  Let \(u,w\) be distinct
variables different from \(x,y,z\).  Each line names the rule, the side
sets \(\Gamma,\Delta\) of the instance and the displayed formulas; every
substitution in it is clean.
\begin{enumerate}
\item[(L1)] \(\chi_0,P(uw)\vdash P(yy),P(uw)\): an axiom, with
      \(\Gamma=\{\chi_0\}\), \(\Delta=\{P(yy)\}\).
\item[(L2)] \(\chi_0,\forall y\,P(uy)\vdash P(yy),P(uw)\): rule
      \((\forall\vdash)\) from (L1), with \(\Gamma=\{\chi_0\}\),
      \(\Delta=\{P(yy),P(uw)\}\), \(\varphi(y)=P(uy)\) and \(t=w\); the
      formula \(P(uy)\) has no bound variable.
\item[(L3)] \(\chi_0\vdash P(yy),P(uw)\): rule \((\forall\vdash)\) from
      (L2), with \(\Gamma=\{\chi_0\}\), the same \(\Delta\),
      \(\varphi(x)=\forall y\,P(xy)\) and \(t=u\), so
      \(\varphi(u)=\forall y\,P(uy)\); the only bound variable of
      \(\varphi(x)\) is \(y\neq u\).  The principal formula \(\chi_0\) is
      already in \(\Gamma\) (\cref{def:p02-w02b-p02-001-g-reading}(1)).
\item[(L4)] \(\chi_0\vdash P(yy),\forall z\,P(uz)\): rule \((\vdash\forall)\)
      from (L3), with \(\Gamma=\{\chi_0\}\), \(\Delta=\{P(yy)\}\), \(v=w\),
      \(\varphi(w)=P(uw)\) and generalized variable \(z\), which does not
      occur in \(P(uw)\); \(w\) is not free in \(\Gamma\cup\Delta\).
\item[(L5)] \(\chi_0\vdash P(yy),\chi_1\): rule \((\vdash\forall)\) from
      (L4), with the same \(\Gamma,\Delta\), \(v=u\),
      \(\varphi(u)=\forall z\,P(uz)\) and generalized variable \(x\), which
      does not occur in \(\varphi(u)\); the substitution \((x/u)\) is clean
      because \(x\) is not bound in \(\varphi(u)\), and gives
      \(\forall z\,P(xz)\); \(u\) is not free in \(\Gamma\cup\Delta\).
\item[(R1)] \(\chi_0,\chi_1,P(yy)\vdash P(yy)\): an axiom, with
      \(\Gamma=\{\chi_0,\chi_1\}\), \(\Delta=\emptyset\).
\item[(R2)] \(\chi_0,\chi_1,\forall z\,P(yz)\vdash P(yy)\): rule
      \((\forall\vdash)\) from (R1), with \(\Gamma=\{\chi_0,\chi_1\}\),
      \(\Delta=\{P(yy)\}\), \(\varphi(z)=P(yz)\) and \(t=y\); \(P(yz)\) has
      no bound variable.
\item[(R3)] \(\chi_0,\chi_1\vdash P(yy)\): rule \((\forall\vdash)\) from
      (R2), with \(\Gamma=\{\chi_0,\chi_1\}\), \(\Delta=\{P(yy)\}\),
      \(\varphi(x)=\forall z\,P(xz)\) and \(t=y\), so
      \(\varphi(y)=\forall z\,P(yz)\); the only bound variable of
      \(\varphi(x)\) is \(z\neq y\).
\item[(C)] \(\chi_0\vdash P(yy)\): cut on \(\chi_1\) from (R3) and (L5),
      with \(\Gamma=\{\chi_0\}\), \(\Delta=\{P(yy)\}\).
\end{enumerate}

\emph{No derivation in \(\mathbf G\).}  Let \(A\) be the set consisting of
\(\chi_0\), the formulas \(\forall y\,P(ty)\) for terms \(t\) with
\(y\notin\operatorname{Var}(t)\), and the atomic formulas \(P(tt')\) for
terms \(t,t'\) with \(y\notin\operatorname{Var}(t)\).  Let \(\mathcal C\)
be the class of sequents \(\Gamma\vdash\Delta\) with \(\Gamma\subseteq A\)
and \(\Delta\subseteq\{P(yy)\}\).  The given sequent lies in
\(\mathcal C\).  We show that no sequent of \(\mathcal C\) has a derivation
in \(\mathbf G\).  Suppose some does, and take a derivation of least height
among all derivations in \(\mathbf G\) of sequents of \(\mathcal C\); let
\(\Gamma\vdash\Delta\) be its end sequent.
\begin{enumerate}
\item[(1)] \(\Gamma\vdash\Delta\) is not an axiom.  The axiom
      \(\Gamma',\varphi\vdash\Delta',\varphi\) needs a formula on both
      sides; the only candidate is \(P(yy)\), which is not in \(A\) because
      its first term contains \(y\).  The equality axiom needs an equation
      on the right, and \(\Delta\) has none.
\item[(2)] So the end sequent is the conclusion of a rule of \(\mathbf G\)
      applied to the sequents immediately above it.  It is not the
      conclusion of a right rule: each right rule puts a formula whose main
      symbol is \(\neg\), \(\land\), \(\lor\), \(\to\), \(\forall\) or
      \(\exists\) into the right side of its conclusion, and \(P(yy)\) is
      atomic.  It is not the conclusion of the equality rule, which puts an
      equation \(E\) into the left side of its conclusion, and \(A\)
      contains no equation.  It is not the conclusion of \((\neg\vdash)\),
      \((\land\vdash)\), \((\lor\vdash)\), \((\to\vdash)\) or
      \((\exists\vdash)\), whose principal formula is a member of the left
      side with main symbol \(\neg\), \(\land\), \(\lor\), \(\to\) or
      \(\exists\); the members of \(A\) are atomic or begin with
      \(\forall\).
\item[(3)] So the last rule is \((\forall\vdash)\).  Its premise is
      \(\Gamma',\varphi(t)\vdash\Delta\) and its conclusion is
      \(\Gamma',\forall v\,\varphi(v)\vdash\Delta\), where
      \(\Gamma'\cup\{\forall v\,\varphi(v)\}=\Gamma\).  So
      \(\Gamma'\subseteq A\) and \(\forall v\,\varphi(v)\in A\), and
      \(v\) and \(\varphi(v)\) are the variable and the scope of the
      leading quantifier.  If \(\forall v\,\varphi(v)=\chi_0\), then
      \(v=x\) and \(\varphi(x)=\forall y\,P(xy)\); the substitution is
      clean, so \(y\notin\operatorname{Var}(t)\), and
      \(\varphi(t)=\forall y\,P(ty)\in A\).  If
      \(\forall v\,\varphi(v)=\forall y\,P(t_0y)\) with
      \(y\notin\operatorname{Var}(t_0)\), then \(v=y\) and
      \(\varphi(y)=P(t_0y)\), and \(\varphi(t)=P(t_0\,t)\), since \(y\) does
      not occur in \(t_0\); this lies in \(A\).  In both cases the premise
      lies in \(\mathcal C\), and its derivation, the subtree above the end
      sequent, has smaller height.
\end{enumerate}
This contradicts the choice of the derivation.  So
\(\chi_0\vdash P(yy)\) has no derivation in \(\mathbf G\).
\end{proof}
\end{annexitem}

\begin{annexitem}{annex:p02-w02b-p02-001-g-variants}{thm:p02-w02b-p02-001-g-variants}{Variants}
\begin{proof}
\emph{Part 1.}  Say that formulas \(\sigma,\sigma'\) are
\emph{equivalent} if \(\mathfrak M\models\sigma[s]\) iff
\(\mathfrak M\models\sigma'[s]\) for every structure \(\mathfrak M\) and
every assignment \(s\) in it.  This relation is reflexive, symmetric and
transitive.

\emph{(a) One renaming step at the top.}  Let \(z\) not occur in
\((Qx\theta)\); we show that \((Qx\theta)\) and \((Qz\,\theta(z/x))\) are
equivalent.  The substitution is clean because \(z\) does not occur in
\(\theta\).  Let \(\mathfrak M\), \(s\) and \(a\in M\) be given.  By
part~4 of \cref{thm:p02-w02b-p02-001-g-substitution} with \(t=z\), whose
value at \(s(\tfrac az)\) is \(a\),
\(\mathfrak M\models\theta(z/x)[s(\tfrac az)]\) iff
\(\mathfrak M\models\theta[s(\tfrac az)(\tfrac ax)]\).  Since \(x\) occurs
in \((Qx\theta)\) and \(z\) does not, \(z\neq x\); so the assignments
\(s(\tfrac az)(\tfrac ax)\) and \(s(\tfrac ax)\) differ at most at \(z\),
and \(z\notin\FV(\theta)\).  By \cref{lem:fol-coincidence},
\(\mathfrak M\models\theta(z/x)[s(\tfrac az)]\) iff
\(\mathfrak M\models\theta[s(\tfrac ax)]\).  Taking this for every
\(a\in M\) (for \(Q=\forall\)) or for some \(a\in M\) (for \(Q=\exists\)),
the satisfaction clauses of \cref{def:p02-w02-p02-005-bw-satisfaction}
give \(\mathfrak M\models(Qz\,\theta(z/x))[s]\) iff
\(\mathfrak M\models(Qx\theta)[s]\).

\emph{(b) Replacement.}  Let \(\sigma,\sigma'\) be equivalent and let
\(\varphi'\) arise from \(\varphi\) by replacing one occurrence of
\(\sigma\) by \(\sigma'\).  Then \(\varphi\) and \(\varphi'\) are
equivalent.  By induction on \(\varphi\): if the occurrence is \(\varphi\)
itself, then \(\varphi=\sigma\) and \(\varphi'=\sigma'\).  Otherwise it lies
in an immediate constituent, so \(\varphi\) is not atomic, and \(\varphi'\)
is \(\neg\alpha'\), \((\alpha'\circ\beta)\), \((\alpha\circ\beta')\) or
\((Qv\alpha')\) for \(\varphi\) equal to \(\neg\alpha\),
\((\alpha\circ\beta)\), \((\alpha\circ\beta)\) or \((Qv\alpha)\), where the
primed constituent arises from the unprimed one by the replacement and is
equivalent to it by the induction hypothesis.  For the connectives the
satisfaction clauses give the claim directly.  For \((Qv\alpha)\), the
induction hypothesis holds at every assignment, in particular at
\(s(\tfrac av)\) for every \(a\in M\), and the clause for \(Q\) gives the
claim.

By (a) and (b), a renaming step, taken forwards or backwards, turns a
formula into an equivalent one; by transitivity so does a finite sequence
of them.  This is the first assertion.  For sequents, let
\(\Gamma'\vdash\Delta'\) be a variant of a valid sequent
\(\Gamma\vdash\Delta\), and let \(s\) in \(\mathfrak M\) satisfy every member
of \(\Gamma'\).  Every member of \(\Gamma\) is a variant of a member of
\(\Gamma'\), so \(s\) satisfies it.  By validity \(s\) satisfies some
\(\delta\in\Delta\); \(\delta\) is a variant of a member of \(\Delta'\),
which \(s\) therefore satisfies.  So \(\Gamma'\vdash\Delta'\) is valid.  As
the variant relation on sequents is symmetric, the converse holds too.

\emph{Part 2.}  We prove by induction on \(\varphi\) the stronger claim:
for every finite set \(W\) of variables there is a variant \(\varphi^{*}\)
of \(\varphi\) with \(\FV(\varphi^{*})=\FV(\varphi)\) whose quantifier
occurrences are on pairwise distinct variables, none of them in \(W\).
\begin{itemize}
\item \(\varphi\) atomic: \(\varphi^{*}=\varphi\), which has no quantifier
      occurrence.
\item \(\varphi=\neg\alpha\): take \(\varphi^{*}=\neg\alpha^{*}\) with
      \(\alpha^{*}\) for \(W\); it is a variant by (V), and its free
      variables and quantifier occurrences are those of \(\alpha^{*}\).
\item \(\varphi=(\alpha\circ\beta)\): take \(\alpha^{*}\) for \(W\),
      \(\beta^{*}\) for \(W\cup\operatorname{QV}(\alpha^{*})\), and
      \(\varphi^{*}=(\alpha^{*}\circ\beta^{*})\).  By (V) used twice,
      \(\varphi\), \((\alpha^{*}\circ\beta)\), \(\varphi^{*}\) are successive
      variants, so \(\varphi^{*}\) is a variant of \(\varphi\) by
      transitivity.  Its free variables are
      \(\FV(\alpha)\cup\FV(\beta)=\FV(\varphi)\).  Its quantifier occurrences
      are those of \(\alpha^{*}\) and of \(\beta^{*}\); within each the
      variables are pairwise distinct, those of \(\beta^{*}\) avoid
      \(\operatorname{QV}(\alpha^{*})\), and none lies in \(W\).
\item \(\varphi=(Qx\theta)\): take \(\theta^{*}\) for \(W\).  By (V),
      \((Qx\theta^{*})\) is a variant of \(\varphi\).  Choose a variable
      \(z\notin W\), \(z\neq x\), that does not occur in \(\theta^{*}\);
      there are infinitely many variables and only finitely many are
      excluded.  Then \(z\) does not occur in \((Qx\theta^{*})\), and one
      renaming step gives \(\varphi^{*}=(Qz\,\theta^{*}(z/x))\), a variant of
      \(\varphi\) by transitivity.  By part~1 of
      \cref{thm:p02-w02b-p02-001-g-substitution} (the substitution is clean)
      the quantifier occurrences of \(\theta^{*}(z/x)\) are those of
      \(\theta^{*}\), on the same variables; these are pairwise distinct and
      outside \(W\), and none is on \(z\), since \(z\) does not occur in
      \(\theta^{*}\).  The outer quantifier is on \(z\notin W\).  For the
      free variables, the same part gives
      \(\FV(\theta^{*}(z/x))=(\FV(\theta^{*})-\{x\})\cup\{z\}\) if
      \(x\in\FV(\theta^{*})\) and \(\FV(\theta^{*})\) otherwise; as
      \(z\notin\FV(\theta^{*})\), in both cases
      \(\FV(\varphi^{*})=\FV(\theta^{*}(z/x))-\{z\}=\FV(\theta^{*})-\{x\}
      =\FV(\theta)-\{x\}=\FV(\varphi)\).
\end{itemize}
In particular \(\FV(\varphi^{*})\subseteq\FV(\varphi)\), as stated.

\emph{Part 3.}  Let \(\Gamma\vdash\Delta\) be a sequent, list the members of
the finite set \(\Gamma\cup\Delta\) as \(\chi_1,\ldots,\chi_k\), and let
\(W=\FV(\chi_1)\cup\cdots\cup\FV(\chi_k)\), a finite set.  By part~2
choose for each \(i\) a variant \(\chi_i^{*}\) of \(\chi_i\) with
\(\FV(\chi_i^{*})=\FV(\chi_i)\), whose quantifier occurrences are on pairwise
distinct variables outside \(W\).  Let
\(\Gamma^{*}=\{\chi_i^{*}:\chi_i\in\Gamma\}\) and
\(\Delta^{*}=\{\chi_i^{*}:\chi_i\in\Delta\}\).  Every member of \(\Gamma\) is
a variant of a member of \(\Gamma^{*}\) and conversely, and likewise for
\(\Delta\) and \(\Delta^{*}\); so \(\Gamma^{*}\vdash\Delta^{*}\) is a variant
of \(\Gamma\vdash\Delta\).  It separates its variables.  (a) A variable
quantified in \(\Gamma^{*}\cup\Delta^{*}\) is quantified in some
\(\chi_i^{*}\) and so lies outside \(W\), while a variable free in some
\(\chi_j^{*}\) lies in \(\FV(\chi_j)\subseteq W\).  (b) In each
\(\chi_i^{*}\) distinct quantifier occurrences are on distinct variables,
so no quantifier on a variable lies in the scope of another quantifier on
the same variable.
\end{proof}
\end{annexitem}

\begin{annexitem}{annex:p02-w02b-p02-001-g-variant-cutfree}{thm:p02-w02b-p02-001-g-variant-cut-elimination}{\(\mathbf G\) and \(\mathbf G^{+}\) on all sequents}
\begin{proof}
Derivations are taken under the renaming reading of
\cref{def:p02-w02b-p02-001-g-variant-reading}.  The variant relation on
sequents is symmetric and transitive, because the variant relation on
formulas is.

\emph{(ii) implies (i).}  Every rule of \(\mathbf G\) is a rule of
\(\mathbf G^{+}\), so a derivation in \(\mathbf G\) under the renaming
reading is one in \(\mathbf G^{+}\).

\emph{(i) implies (iii).}  Every axiom of \(\mathbf G^{+}\) is valid, and
every instance of a rule of \(\mathbf G^{+}\), read as in
\cref{def:p02-w02b-p02-001-g-reading}, with valid premises has a valid
conclusion; this is the soundness of \(\mathbf G^{+}\) established in the
proof of \cref{thm:p02-w02-p02-005-bw-completeness-g} and used in the first
paragraph of the proof of \cref{thm:p02-w02b-p02-001-g-completeness}.  We
show by induction on the height of the subtree above a node that every
node of a derivation in \(\mathbf G^{+}\) under the renaming reading is
valid.  If the node is a variant of an axiom, it is valid by part~1 of
\cref{thm:p02-w02b-p02-001-g-variants}.  Otherwise there is a rule
instance whose conclusion is a variant of the node and whose premises are
variants of the nodes immediately above it, one for each.  Those nodes are
valid by the induction hypothesis, so the premises are valid by part~1 of
\cref{thm:p02-w02b-p02-001-g-variants}, so the conclusion is valid by
soundness, and so is the node, again by part~1.  The end sequent is the
lowest node.

\emph{(iii) implies (ii).}  Let \(\Gamma\vdash\Delta\) be valid.  By part~3
of \cref{thm:p02-w02b-p02-001-g-variants} it has a variant
\(\Gamma'\vdash\Delta'\) that separates its variables, and
\(\Gamma'\vdash\Delta'\) is valid by part~1.  By
\cref{thm:p02-w02b-p02-001-g-completeness} it has a derivation \(D\) in
\(\mathbf G\) in the sense of \cref{def:p02-w02b-p02-001-g-reading}: each
node is an axiom or is the conclusion of a rule instance of \(\mathbf G\)
whose premises are the nodes immediately above it.  Every node of \(D\)
therefore meets the condition of the renaming reading, each sequent being a
variant of itself.  Let \(D'\) be \(D\) with the label of the lowest node
changed from \(\Gamma'\vdash\Delta'\) to \(\Gamma\vdash\Delta\).  The nodes
of \(D'\) other than the lowest are those of \(D\), with the same nodes
above them.  The lowest node of \(D'\) is a variant of
\(\Gamma'\vdash\Delta'\), which is an axiom or the conclusion of a rule
instance of \(\mathbf G\) whose premises are the nodes immediately above;
so it meets the condition too.  Thus \(D'\) is a derivation of
\(\Gamma\vdash\Delta\) in \(\mathbf G\) under the renaming reading.

\emph{The consequences.}  By (i) iff (ii), \(\mathbf G^{+}\) and
\(\mathbf G\) have the same derivable sequents under the renaming reading.
Let \(T\) be a set of sentences and \(\varphi\) a sentence.  If
\(T\models\varphi\), then by \cref{thm:p02-w02-p02-005-bw-completeness-g}
some sequent \(\Gamma\vdash\{\varphi\}\) with finite \(\Gamma\subseteq T\)
has a derivation in \(\mathbf G^{+}\); a derivation in the sense of
\cref{def:p02-w02b-p02-001-g-reading} is one under the renaming reading,
so by (i) implies (ii) the sequent has a derivation in \(\mathbf G\).
Conversely, if \(\Gamma\vdash\{\varphi\}\) with finite \(\Gamma\subseteq T\)
has a derivation in \(\mathbf G\), it is valid by (ii) implies (iii), so
\(\Gamma\models\varphi\) by part~1 of
\cref{def:p02-w02b-p02-001-g-valid}; every model of \(T\) is a model of
\(\Gamma\), so \(T\models\varphi\).
\end{proof}
\end{annexitem}

\section{Proofs for Lindstr\"om's theorem in Barwise's formulation}
\label{sec:annex-p02-w02b-p02-001-lindstrom}

The notation is that of \cref{sec:p02-w02b-p02-001-lindstrom}.

\begin{annexitem}{annex:p02-w02b-p02-001-barwise-lindstrom}{thm:p02-w02b-p02-001-barwise-lindstrom}{Barwise's 5.7 from the proved form}
\begin{proof}
Let \(\mathcal L^{\ast}\) satisfy the hypotheses.  We verify the two
numbered hypotheses of \cref{thm:p03-lindstrom}.

\emph{Hypothesis~1.}  Let \(T\) be a countable set of
\(\mathcal L^{\ast}\)-sentences every finite subset of which is
satisfiable.  Let \(\tau\) be the union of the finite occurrence signatures
of the members of \(T\); it is a relational signature, and the inclusion of
each occurrence signature into \(\tau\) is an arity-preserving renaming, so
by the renaming clause of regularity each member of \(T\) is a sentence of
\(\tau\) with the same satisfaction relation.  Let \(T_0\subseteq T\) be
finite, with a model \(\mathfrak M_0\).  Take the reduct of
\(\mathfrak M_0\) to the union of the occurrence signatures of the members
of \(T_0\) and expand it to \(\tau\) by interpreting every further
relation symbol as the empty relation.  Satisfaction of each member of
\(T_0\) depends only on the reduct to its occurrence signature, which is
unchanged, so the resulting \(\tau\)-structure is a model of \(T_0\).  Thus
every finite \(T_0\subseteq T\) has a \(\tau\)-model, and part~1 of
\cref{def:p02-w02b-p02-001-lindstrom-reading}, applied to \(\tau\) and
\(T\), gives a model of \(T\).  No use is made of the countability of
\(T\).

\emph{Hypothesis~2.}  Let \(\alpha\) be a satisfiable
\(\mathcal L^{\ast}\)-sentence with occurrence signature \(\tau\).  The set
\(\{\alpha\}\) has one member, so it is a set of at most \(\aleph_0\)
sentences of \(\tau\) that has a model.  Part~2 of
\cref{def:p02-w02b-p02-001-lindstrom-reading} gives a model of \(\alpha\)
whose domain has cardinality at most \(\aleph_0\), that is, a countable
model.  The proof of \cref{thm:p03-lindstrom} uses this hypothesis to
enumerate the two sorts of such a model, which is possible for a finite as
well as for a countably infinite domain.

Both hypotheses of \cref{thm:p03-lindstrom} hold, so every
\(\mathcal L^{\ast}\)-sentence is equivalent, over its finite occurrence
signature, to a relational first-order sentence.  This is the conclusion,
and it is part~3 of \cref{def:p02-w02b-p02-001-lindstrom-reading}.
\end{proof}
\end{annexitem}

\begin{annexitem}{annex:p02-w02b-p02-001-lindstrom-sharp}{thm:p02-w02b-p02-001-lindstrom-sharp}{The two hypotheses for \(L(Q)\)}
\begin{proof}
\emph{Part~1.}  Suppose \(Qx\,(x=x)\) were equivalent to a first-order
sentence \(\sigma\) of the empty signature, so that a structure, which here
is a nonempty set, satisfies \(\sigma\) iff its domain is uncountable.  The
set \(\mathbb R\) of reals is uncountable, so it is a model of
\(\{\sigma\}\).  By \cref{thm:p02-w02-p02-005-bw-ls} with
\(\kappa=\aleph_0\), applied to the one-element set \(\{\sigma\}\), there is
a model of \(\sigma\) whose domain has cardinality at most \(\aleph_0\).
That domain is countable, so it does not satisfy \(Qx\,(x=x)\), contrary to
the assumed equivalence.

\emph{Part~2, finite subsets.}  Let \(T_0\subseteq T\) be finite and let
\(\alpha_1,\ldots,\alpha_n\) be the distinct ordinals \(\alpha\) such that
\(c_\alpha\) occurs in a sentence of \(T_0\); if there are none, put
\(n=1\) and choose any \(\alpha_1<\omega_1\).  Let \(\mathfrak M\) have
domain \(M=\{1,\ldots,n\}\), \(P^{\mathfrak M}=M\),
\(c_{\alpha_i}^{\mathfrak M}=i\) for \(i\leq n\), and
\(c_\beta^{\mathfrak M}=1\) for every other \(\beta<\omega_1\).  The set
\(P^{\mathfrak M}\) is finite, so \(\mathfrak M\models\neg Qx\,P(x)\).
Every member of \(T_0\) of the form \(P(c_\alpha)\) has
\(\alpha=\alpha_i\) for some \(i\), and \(i\in P^{\mathfrak M}\).  Every
member of the form \(\neg\,c_\alpha=c_\beta\) with \(\alpha<\beta\) has
\(\alpha=\alpha_i\) and \(\beta=\alpha_j\) with \(i\ne j\), because the
\(\alpha_k\) are distinct, and \(i\ne j\) in \(M\).  So
\(\mathfrak M\) is a model of \(T_0\).

\emph{Part~2, the whole set.}  Suppose \(\mathfrak M\models T\).  The map
\(\alpha\mapsto c_\alpha^{\mathfrak M}\) on \(\omega_1\) takes values in
\(P^{\mathfrak M}\), because \(P(c_\alpha)\in T\), and is one-to-one,
because \(\neg\,c_\alpha=c_\beta\in T\) for \(\alpha<\beta\).  Since
\(\omega_1\) is uncountable, so is \(P^{\mathfrak M}\), and therefore
\(\mathfrak M\models Qx\,P(x)\), contradicting
\(\neg Qx\,P(x)\in T\).  So \(T\) has no model, and 2.4, read for
\(L(Q)\)-sentences of this language, fails for \(T\).
\end{proof}
\end{annexitem}
